
\documentclass{article}
\usepackage{colortbl}
\usepackage{makecell}
\usepackage{multirow}
\usepackage{supertabular}

\begin{document}

\newcounter{utterance}

\centering \large Interaction Transcript for game `cladder', experiment `full\_v1.5\_default', episode 9980 with refineonly{-}Qwen3.5{-}27B.
\vspace{24pt}

{ \footnotesize  \setcounter{utterance}{1}
\setlength{\tabcolsep}{0pt}
\begin{supertabular}{c@{$\;$}|p{.15\linewidth}@{}p{.15\linewidth}p{.15\linewidth}p{.15\linewidth}p{.15\linewidth}p{.15\linewidth}}
    \# & \multicolumn{2}{c}{Player} && \multicolumn{2}{c}{Game Master} \\
    \hline

    \theutterance \stepcounter{utterance}  
    & & & \multicolumn{4}{p{0.6\linewidth}}{
        \cellcolor[rgb]{0.9,0.9,0.9}{
            \makecell[{{p{\linewidth}}}]{
                \texttt{\tiny{[P1$\langle$GM]}}
                \texttt{You are an expert in causal inference. The following question is not a typical commonsense query, but rather a meticulously designed question created by a professor specializing in causal inference, intended to assess the students' mastery of the course content.} \\
\texttt{Imagine a self{-}contained, hypothetical world with only the following conditions, and without any unmentioned factors or causal relationships: Hwax has a direct effect on pexu and rukz. Pexu has a direct effect on rukz. We know that hwax causes pexu. hwax and pexu causes rukz. We observed an individual is not hwax. Would an individual is not rukz if not pexu instead of pexu?} \\
\texttt{Start your answer with 'Yes' or 'No', followed by additional reasoning or evidence to support your explanation.} \\
            }
        }
    }
     \\ \\

    \theutterance \stepcounter{utterance}  
    & \multicolumn{4}{p{0.6\linewidth}}{
        \cellcolor[rgb]{0.561,0.737,0.561}{
            \makecell[{{p{\linewidth}}}]{
                \texttt{\tiny{[P1$\rangle$GM]}}
                \texttt{No} \\
\\ 
\texttt{**Reasoning:**} \\
\\ 
\texttt{To determine the answer, we must reconstruct the causal structure and evaluate the counterfactual scenario based strictly on the provided conditions.} \\
\\ 
\texttt{1.  **Causal Structure**:} \\
\texttt{    *   \$Hwax \textbackslash{}to Pexu\$} \\
\texttt{    *   \$Hwax \textbackslash{}to Rukz\$} \\
\texttt{    *   \$Pexu \textbackslash{}to Rukz\$} \\
\texttt{    *   This forms a Directed Acyclic Graph (DAG) where \$Hwax\$ is a common cause of \$Pexu\$ and \$Rukz\$, and \$Pexu\$ is also a direct cause of \$Rukz\$.} \\
\\ 
\texttt{2.  **Observed Facts**:} \\
\texttt{    *   The individual is **not** \$Hwax\$ (\$Hwax = 0\$).} \\
\texttt{    *   The individual **is** \$Pexu\$ (\$Pexu = 1\$).} \\
\texttt{    *   The individual **is** \$Rukz\$ (\$Rukz = 1\$). (This is implied by the counterfactual phrasing "if not pexu instead of pexu," which presupposes the current state is \$Pexu\$ and the outcome is \$Rukz\$).} \\
\\ 
\texttt{3.  **The Counterfactual Question**:} \\
\texttt{    *   "Would an individual be not \$Rukz\$ if not \$Pexu\$ instead of \$Pexu\$?"} \\
\texttt{    *   Formally: Given \$Hwax = 0\$ and \$Pexu = 1\$, what is the value of \$Rukz\$ in the counterfactual world where we intervene to set \$Pexu = 0\$? (i.e., Calculate \$Rukz\_\{Pexu=0\}\$).} \\
\\ 
\texttt{4.  **Analysis of the Counterfactual**:} \\
\texttt{    *   In the counterfactual world, we intervene to set \$Pexu = 0\$.} \\
\texttt{    *   The variable \$Hwax\$ remains at its observed value of \$0\$ (since the intervention is only on \$Pexu\$).} \\
\texttt{    *   We must determine the value of \$Rukz\$ given \$Hwax = 0\$ and \$Pexu = 0\$.} \\
\texttt{    *   The problem states that \$Hwax\$ has a direct effect on \$Rukz\$ and \$Pexu\$ has a direct effect on \$Rukz\$. However, the problem **does not specify the functional form** (the mechanism) of these effects.} \\
\texttt{        *   Does \$Hwax\$ alone cause \$Rukz\$? (If \$Hwax=0\$, does \$Rukz\$ become 0?)} \\
\texttt{        *   Does \$Pexu\$ alone cause \$Rukz\$?} \\
\texttt{        *   Is \$Rukz\$ an} \\
            }
        }
    }
    & & \\ \\

    \theutterance \stepcounter{utterance}  
    & & & \multicolumn{2}{p{0.3\linewidth}}{
        \cellcolor[rgb]{0.9,0.9,0.9}{
            \makecell[{{p{\linewidth}}}]{
                \texttt{\tiny{[GM$|$GM]}}
                \texttt{No} \\
            }
        }
    }
    & & \\ \\

    \theutterance \stepcounter{utterance}  
    & & & \multicolumn{2}{p{0.3\linewidth}}{
        \cellcolor[rgb]{0.9,0.9,0.9}{
            \makecell[{{p{\linewidth}}}]{
                \texttt{\tiny{[GM$|$GM]}}
                \texttt{yes} \\
            }
        }
    }
    & & \\ \\

    \theutterance \stepcounter{utterance}  
    & & & \multicolumn{2}{p{0.3\linewidth}}{
        \cellcolor[rgb]{0.9,0.9,0.9}{
            \makecell[{{p{\linewidth}}}]{
                \texttt{\tiny{[GM$|$GM]}}
                \texttt{game\_result = LOSE} \\
            }
        }
    }
    & & \\ \\

\end{supertabular}
}

\end{document}
